\section{Conclusion}

Functional Loops capture a failure mode that tool-name repetition misses:
after feedback is visible, tool-using LLMs may issue different calls that make no
progress and repeat the same unresolved functional effect. This matters because
the model continues spending interaction budget without changing the state
that blocks completion. The pattern depends on both the model and the
environment, and a Functional Loop can still precede a useful retry. On STATE-Bench,
however, Functional Loop-exposed trajectories are associated with lower scores and more
calls. A Functional Loop is therefore a useful recovery signal, but not a reason to
terminate every loop.

The experiments show a consistent task-level benefit. Every reported
Raw/PERSIST-ACE pair improves its primary outcome across STATE-Bench,
DeepPlanning Shopping, and $\tau$-Knowledge. The trend persists when different
models are evaluated on the same benchmark and when the same model is evaluated
on different benchmarks. Better outcomes do not require uniformly longer successful runs:
calls per successful trajectory usually decrease and otherwise remain close.
The ablations further locate the advantage in the recovery design. Online
Replanning with a generic prompt does not recover the baseline, the Validated Registry intervenes
more selectively and with lower observed loss than the Unvalidated All-Card
Registry, and complete Cards outperform operator-only and ungrounded variants
in held-out same-state comparisons. Independently removing the soft controller
gate increases coverage but not task performance.

Taken together, the results show why a Functional Loop is a practical unit for LLM tool-use
recovery. No progress recalls a candidate; repeated effect confirms the Functional Loop;
Card admission then decides whether it is recoverable. Across heterogeneous
models and tool environments, this separation improves task outcomes while
protecting productive retries and avoiding systematic growth in
successful-trajectory interaction cost.
